% ------------------------------------------------------------
\escenario{ESC-01}{Colocar una barrera con el tablero avanzado}

\begin{tabla}{|L{3.2cm}|Y|}
\fichacab
\bloque{Trazabilidad}
\parte{Característica}{QA-01 Desempeño: comportamiento temporal.}
\parte{Stakeholders}{STK-01 Jugador (8 años o más), STK-09 Diseñador de experiencia de usuario.}
\parte{Concerns}{CRN-02: \emph{``Cuando coloco una barrera quiero que el juego me conteste al instante. Si tarda más de un segundo en decirme si vale, la partida se me hace pesada.''}}
\parte{Restricciones y dependencias}{REQ-02 (respuesta menor a 1 segundo), REQ-01 (multijugador en línea obligatorio), CON-02 (TCP), DEP-04 (infraestructura externa).}
\parte{Tensiones y preguntas}{TEN-01 Latencia contra integridad. TEN-06 Trazabilidad contra tiempo de respuesta.}
\parte{Tipo y prioridad}{De uso, en tiempo de ejecución. Prioridad (A, A).}
\bloque{Escenario}
\parte{Fuente del estímulo}{Jugador en su turno durante una partida en línea.}
\parte{Estímulo}{Coloca una barrera en el tablero y confirma la jugada.}
\parte{Ambiente}{Operación normal en el peor caso realista: tablero de 9 × 9 con 16 de las 20 barreras ya colocadas, ambos jugadores conectados desde redes domésticas con una latencia de ida y vuelta de hasta 150 ms, y el servidor atendiendo otras 20 partidas al mismo tiempo.}
\parte{Artefacto}{Motor de reglas (incluida la validación de que ambos jugadores conservan un camino a su meta), servidor de partidas, gestión de turnos y estado, y capa de red sobre TCP.}
\parte{Respuesta}{El jugador ve si su barrera quedó aceptada o rechazada, y el oponente ve exactamente el mismo tablero. Una barrera que se le mostró como aceptada no desaparece después.}
\parte{Medida de la respuesta}{En 1\,000 colocaciones medidas: el 100~\% recibe respuesta en menos de 1 segundo, contado desde que el jugador confirma hasta que ve el resultado; la mediana queda por debajo de 300 ms; el oponente ve la barrera en menos de 1 segundo en el 100~\% de los casos; y 0 barreras mostradas como aceptadas se revierten con un cliente no alterado.}
\end{tabla}

\textbf{Por qué es relevante.} Colocar una barrera es la acción más frecuente del juego junto con mover la ficha, y es la más costosa: es la única que obliga a buscar un camino para los dos jugadores antes de aceptarla. Como el modo en línea es obligatorio, esa búsqueda no puede quedarse en el equipo del jugador: tiene que pasar por el servidor, y en ese mismo segundo compiten el viaje por TCP, la validación autoritativa que pide seguridad (TEN-01) y la escritura que piden auditoría y recuperación (TEN-06). Es el escenario donde se ve si REQ-02 se puede cumplir con las demás características en pie. Se eligió el tablero avanzado porque es cuando la búsqueda de camino tiene más trabajo: los caminos son más largos y, cuando una barrera encierra al rival, hay que recorrer toda la zona alcanzable para demostrarlo.

\textbf{Cómo se verifica.} Un cliente de prueba automatizado reproduce posiciones reales de final de partida y registra la hora de confirmación y la de respuesta en cada colocación. La latencia de 150 ms se simula con un emulador de red, y las 20 partidas simultáneas con clientes automáticos. Se reportan la mediana, el máximo y el número de barreras revertidas.
